\documentclass[11pt,a4paper]{article}

%Inserindo pacotes básicos%

\usepackage[utf8]{inputenc} 
\usepackage[T1]{fontenc} 
\usepackage[brazil]{babel} 
\usepackage{mathptmx} 
\usepackage{geometry} 
\usepackage{setspace} 
\usepackage{indentfirst}
\usepackage{graphicx} 
\usepackage{float} 
\usepackage{amsmath} 
\usepackage{siunitx} 
\usepackage{booktabs}
\usepackage{array} 
\usepackage{caption} 
\usepackage{fancyhdr}
\setlength{\headheight}{13.6pt}
\usepackage{ragged2e} 
\usepackage{url} 
\usepackage{hyperref}

%Configurações gerais%

\geometry{ 
    a4paper, 
    top=3cm, 
    bottom=2cm, 
    left=3cm, 
    right=2cm 
} 
\onehalfspacing 
\justifying 
\setlength{\parindent}{1.25cm} 
\setlength{\parskip}{0pt} 
\sisetup{ output-decimal-marker = {,}, per-mode = symbol 
}

%Cabeçalho e rodapé%

\pagestyle{fancy} 
\fancyhf{} 
\fancyhead[L]{Verificação da Relutância de Circuito Magnético} 
\fancyhead[R]{CEME I} 
\fancyfoot[C]{\thepage} 
\renewcommand{\headrulewidth}{0.4pt}

%Comando auxiliar para figuras%

\newcommand{\figuraopcional}[4]{ 
    \IfFileExists{#1}{ 
        \begin{figure}[H] 
            \centering 
            \includegraphics[width=#2\textwidth]{#1} 
            \caption{#3} 
            \label{#4} 
        \end{figure} 
    }{ 
        \begin{figure}[H] 
            \centering 
            \fbox{ 
                \begin{minipage}[c][4cm][c]{0.85\textwidth} 
                \centering 
                Figura não carregada no Overleaf.\\ 
                Arquivo esperado: \texttt{#1} 
            \end{minipage} 
        } 
        \caption{#3} 
        \label{#4} 
    \end{figure} 
}
}
%Dados do relatorio%
\title{ 
    \textbf{Verificação da Relutância de Circuito Magnético}\\ 
    \large Universidade do Estado do Rio de Janeiro\\ 
    \large Conversão Eletromecânica de Energia I 
} 
    
\author{ 
    Turma 2\\ 
    João, Elton e Emanuel 
} 
\date{14 de agosto de 2026} 
\begin{document}

\maketitle 
\thispagestyle{fancy}

\begin{abstract} Este relatório apresenta o estudo experimental da relutância em circuitos magnéticos utilizando núcleos ferromagnéticos nas configurações U+I e E+I. Inicialmente foram realizadas medições com o núcleo U+I aberto e fechado, observando-se a influência do fechamento do caminho magnético sobre a corrente de magnetização e sobre a tensão induzida em uma bobina de teste.

Em seguida, foi analisada a distribuição do fluxo magnético em um núcleo E+I, considerando a bobina de alimentação posicionada em um ramo lateral e posteriormente no ramo central. Os fluxos magnéticos foram calculados indiretamente por meio da tensão induzida na bobina de teste, utilizando a lei de Faraday.

Também foi realizada uma simulação no software FEMM, com material magnético linear de permeabilidade relativa igual a 1000, a fim de comparar qualitativamente a distribuição de fluxo obtida experimentalmente com a distribuição simulada. Os resultados mostraram que o fechamento do circuito magnético reduz significativamente a relutância e aumenta o fluxo produzido, enquanto o núcleo E+I apresentou divisão de fluxo coerente com a lei dos nós aplicada a circuitos magnéticos. \end{abstract}

\section{Introdução} Circuitos magnéticos são estruturas formadas por materiais de elevada permeabilidade magnética, geralmente utilizadas para conduzir e concentrar fluxo magnético em máquinas elétricas, transformadores, atuadores e outros dispositivos eletromecânicos. De forma análoga aos circuitos elétricos, nos quais a corrente é estabelecida pela aplicação de uma força eletromotriz sobre uma determinada resistência, em circuitos magnéticos o fluxo magnético é produzido por uma força magnetomotriz aplicada sobre uma determinada relutância.

A força magnetomotriz é dada por 
\begin{equation} 
    \mathcal{F} = N I , 
    \label{eq:fmm} 
\end{equation}

em que \(N\) representa o número de espiras da bobina e \(I\) representa a corrente que circula por ela. A relutância magnética, por sua vez, pode ser entendida como a oposição oferecida pelo circuito à passagem do fluxo magnético. Para um caminho magnético uniforme, a relutância pode ser expressa por

\begin{equation} 
    \mathcal{R} = \frac{l}{\mu A}, 
    \label{eq:relutancia} 
\end{equation}

em que \(l\) é o comprimento médio do caminho magnético, \(A\) é a área da seção transversal e \(\mu\) é a permeabilidade magnética do material. A permeabilidade absoluta pode ser escrita como

\begin{equation} 
    \mu = \mu_0 \mu_r , 
    \label{eq:permeabilidade} 
\end{equation}

sendo \(\mu_0\) a permeabilidade magnética do vácuo e \(\mu_r\) a permeabilidade relativa do material.

A relação entre força magnetomotriz, relutância e fluxo magnético é apresentada por

\begin{equation} 
    \Phi = \frac{\mathcal{F}}{\mathcal{R}} = \frac{N I}{\mathcal{R}}, \label{eq:fluxo_relutancia} 
\end{equation}

onde \(\Phi\) é o fluxo magnético no circuito. Essa relação mostra que, para uma mesma força magnetomotriz, a redução da relutância resulta em maior fluxo magnético. Por esse motivo, a presença de entreferros ou descontinuidades no caminho magnético aumenta consideravelmente a relutância do circuito e reduz o fluxo produzido.

Como a medição direta do fluxo magnético em um núcleo ferromagnético é difícil, o experimento utiliza uma bobina de teste para realizar uma medição indireta. De acordo com a lei de Faraday, a tensão induzida em uma bobina é proporcional à taxa de variação do fluxo magnético concatenado com suas espiras. Para uma tensão senoidal, o valor eficaz do fluxo magnético pode ser calculado por

\begin{equation} 
    \Phi_{\text{rms}} = \frac{V_{\text{rms}}}{N \omega}, 
    \label{eq:fluxo_farad} 
\end{equation}

em que \(V_{\text{rms}}\) é a tensão eficaz induzida na bobina de teste, \(N\) é o número de espiras e \(\omega\) é a frequência angular da fonte, dada por

\begin{equation} 
    \omega = 2 \pi f . 
    \label{eq:omega} 
\end{equation}

Neste experimento, a frequência utilizada foi de \(\SI{60}{Hz}\), resultando em

\begin{equation} 
    \omega = 2 \pi \cdot 60 \approx \SI{376,99}{rad/s}. 
    \label{eq:omega_valor} 
\end{equation}

A análise dos circuitos magnéticos também permite aplicar uma forma análoga à lei dos nós. Em um ponto de divisão de fluxo, o fluxo que entra em uma junção deve ser aproximadamente igual à soma dos fluxos que saem, desde que sejam desconsideradas perdas por dispersão. Assim, para o núcleo E+I, espera-se que o fluxo no ramo de alimentação se divida entre os demais ramos, respeitando aproximadamente a relação

\begin{equation}
    \Phi_{\text{entrada}} \approx \Phi_{\text{ramo 1}} + \Phi_{\text{ramo 2}} . \label{eq:lei_nos} 
\end{equation}

Essa aproximação é válida dentro das limitações experimentais, pois parte do fluxo pode se dispersar pelo ar, além de existirem imperfeições de montagem, folgas mecânicas, não linearidades do material e incertezas associadas às medições.

\section{Objetivo} O objetivo deste relatório é verificar experimentalmente a influência da relutância no comportamento de circuitos magnéticos, por meio de ensaios realizados com núcleos ferromagnéticos nas configurações U+I e E+I. Busca-se calcular os fluxos magnéticos a partir das tensões induzidas em uma bobina de teste, comparar os resultados obtidos para diferentes configurações do núcleo, verificar a validade aproximada da lei dos nós magnética e comparar os resultados experimentais com uma simulação realizada no software FEMM.

\section{Procedimentos} 
\subsection{Materiais utilizados} 
Foram utilizados os seguintes materiais e equipamentos: 
\begin{itemize} 
    \item Duas bobinas de \num{1000} espiras; 
    \item Fonte de tensão alternada; 
    \item Voltímetro CA; 
    \item Amperímetro CA; 
    \item Núcleo ferromagnético em formato U; 
    \item Núcleo ferromagnético em formato I pequeno;
    \item Núcleo ferromagnético em formato E; 
    \item Núcleo ferromagnético em formato I grande; 
    \item Software FEMM para simulação do circuito magnético. 
\end{itemize}
\subsection{Dimensões dos núcleos}

\begin{table}[H] 
    \centering 
    \caption{Dimensões dos núcleos utilizados no experimento.} \label{tab:dimensoes_nucleos} 
    \begin{tabular}{lccc} 
        \toprule Núcleo & Altura (cm) & Largura (cm) & Profundidade (cm) \\ 
        \midrule Núcleo pequeno & 17 & 15 & 4 \\ Núcleo grande & 25 & 25 & 4 \\ \bottomrule 
    \end{tabular} 
\end{table}

A distância entre os núcleos foi medida como aproximadamente \(\SI{6,5}{cm}\). Para a modelagem no FEMM, foi adotada a profundidade de \(\SI{4}{cm}\), compatível com a medição experimental.

\subsection{Ensaio com núcleo U+I aberto}

Na primeira etapa, foi montado o circuito magnético utilizando o núcleo em formato U, a fonte de tensão alternada e duas bobinas. Uma das bobinas foi utilizada como bobina de alimentação e a outra como bobina de teste. A fonte foi ajustada para alimentar a bobina principal com tensão de \(\SI{15}{V}\) em corrente alternada. Em seguida, foram medidos a corrente na bobina de alimentação e a tensão induzida na bobina de teste.

\subsection{Ensaio com núcleo U+I fechado}

Na segunda etapa, foi adicionada a peça em formato I ao núcleo U, fechando o caminho magnético. A bobina de alimentação foi novamente alimentada com \(\SI{15}{V}\) em corrente alternada. Foram medidos a corrente na bobina de alimentação e a tensão induzida na bobina de teste. Esse procedimento permitiu comparar diretamente o comportamento do circuito magnético aberto e fechado.

\subsection{Ensaio com núcleo E+I fechado}

Na terceira etapa, foi utilizado o núcleo maior na configuração E+I. Inicialmente, a bobina de alimentação foi posicionada no ramo lateral esquerdo do núcleo, enquanto a bobina de teste foi colocada sucessivamente nas posições 1, 2 e 3. Para cada posição, foi medida a tensão induzida na bobina de teste.

Posteriormente, a bobina de alimentação foi posicionada no ramo central do núcleo. O mesmo procedimento foi repetido, medindo-se a tensão induzida nas três posições da bobina de teste. Em todas as medições, a fonte foi ajustada para aproximadamente \(\SI{15}{V}\) em corrente alternada.

\subsection{Simulação no FEMM}

A simulação foi realizada no software FEMM, considerando um problema magnético planar com profundidade de \(\SI{4}{cm}\). O núcleo E+I foi modelado com as dimensões medidas em laboratório. O material magnético foi definido com curva B-H linear e permeabilidade relativa

\begin{equation} 
    \mu_r = 1000 . 
    \label{eq:mur_femm} 
\end{equation}

A bobina de alimentação foi configurada com \num{1000} espiras e corrente de \(\SI{0,014}{A}\), valor correspondente à corrente medida experimentalmente no ensaio com núcleo fechado. Dessa forma, a força magnetomotriz considerada na simulação foi

\begin{equation}    
    \mathcal{F}_{\text{FEMM}} = 1000 \cdot 0,014 = \SI{14}{A.espiras}. \label{eq:fmm_femm} 
\end{equation}

A simulação foi utilizada para observar a distribuição da densidade de fluxo magnético no núcleo e comparar qualitativamente o comportamento obtido com os resultados experimentais.

\section{Resultados e Discussão}

\subsection{Cálculo dos fluxos magnéticos}

Como as bobinas utilizadas possuem \num{1000} espiras e a frequência da fonte é de \(\SI{60}{Hz}\), o fluxo magnético eficaz foi calculado por meio de \eqref{eq:fluxo_farad}:

\begin{equation} 
    \Phi_{\text{rms}} = 
    \frac{V_{\text{rms}}}{1000 \cdot 376,99}. 
    \label{eq:fluxo_calculo} 
\end{equation}

Essa expressão foi utilizada para converter as tensões medidas na bobina de teste em valores aproximados de fluxo magnético.

\subsection{Ensaio com núcleo U+I aberto e fechado}

Os resultados medidos nos ensaios com núcleo U+I aberto e fechado são apresentados na Tabela \ref{tab:ui_resultados}.

\begin{table}[H] 
    \centering 
    \caption{Resultados experimentais dos ensaios com núcleo U+I.} \label{tab:ui_resultados} 
    \begin{tabular}{lccc} 
        \toprule Condição do núcleo & \(V_{\text{entrada}}\) (V) & \(I_{\text{entrada}}\) (A) & \(V_{\text{teste}}\) (V) \\ 
        \midrule 
        U+I aberto & 15 & 0,120 & 4 \\ 
        U+I fechado & 15 & 0,014 & 12 \\ 
        \bottomrule 
    \end{tabular} 
\end{table}

Para o núcleo aberto, o fluxo magnético calculado foi

\begin{equation}
    \Phi_{\text{aberto}} = 
    \frac{4}{1000 \cdot 376,99} = 1,06 \times 10^{-5} \; \text{Wb}. 
    \label{eq:fluxo_aberto} 
\end{equation}

Para o núcleo fechado, o fluxo magnético calculado foi

\begin{equation} 
    \Phi_{\text{fechado}} = 
    \frac{12}{1000 \cdot 376,99} = 3,18 \times 10^{-5} \; 
    \text{Wb}. 
    \label{eq:fluxo_fechado} 
\end{equation}

A comparação entre os dois casos mostra que o fechamento do núcleo aumentou o fluxo magnético em aproximadamente três vezes, uma vez que a tensão induzida na bobina de teste passou de aproximadamente \(\SI{4}{V}\) para \(\SI{12}{V}\). Esse comportamento é coerente com a teoria de circuitos magnéticos, pois a inserção da peça I reduz a relutância do circuito ao fechar o caminho preferencial para o fluxo.

Além do aumento do fluxo, observou-se uma redução expressiva da corrente de alimentação. No núcleo aberto, a corrente medida foi de \(\SI{0,120}{A}\), enquanto no núcleo fechado a corrente foi de apenas \(\SI{0,014}{A}\). Essa redução indica que, com menor relutância, o circuito magnético exige menor corrente de magnetização para estabelecer fluxo no núcleo.

Uma estimativa da relutância equivalente pode ser obtida por

\begin{equation} 
    \mathcal{R} = \frac{N I}{\Phi}. 
    \label{eq:relutancia_exp} 
\end{equation}

Para o núcleo aberto,

\begin{equation} 
    \mathcal{R}_{\text{aberto}} = 
    \frac{1000 \cdot 0,120}{1,06 \times 10^{-5}} = 
    1,13 \times 10^{7} 
    \; \text{A.espiras/Wb}. 
    \label{eq:relutancia_aberto} 
\end{equation}

Para o núcleo fechado, 

\begin{equation} 
    \mathcal{R}_{\text{fechado}} = 
    \frac{1000 \cdot 0,014}{3,18 \times 10^{-5}} =
    4,40 \times 10^{5} 
    \; \text{A.espiras/Wb}. 
    \label{eq:relutancia_fechado} 
\end{equation}

A relutância estimada no caso aberto foi muito maior que a relutância estimada no caso fechado. Esse resultado confirma que o entreferro e a abertura do caminho magnético aumentam significativamente a oposição à passagem do fluxo.

\subsection{Ensaio E+I com alimentação no ramo lateral}

No ensaio com o núcleo E+I, inicialmente a bobina de alimentação foi posicionada no ramo lateral esquerdo. As tensões medidas na bobina de teste e os fluxos calculados são apresentados na Tabela \ref{tab:ei_lateral}.

\begin{table}[H] 
    \centering 
    \caption{Resultados experimentais do núcleo E+I com alimentação no ramo lateral.} \label{tab:ei_lateral} 
    \begin{tabular}{ccc} 
        \toprule Posição da bobina de teste & Tensão medida (V) & Fluxo calculado (Wb) \\ 
        \midrule 
        1 & 14 & \(3,71 \times 10^{-5}\) \\ 
        2 & 9 & \(2,39 \times 10^{-5}\) \\ 
        3 & 3 & \(7,96 \times 10^{-6}\) \\ 
        \bottomrule 
    \end{tabular} 
\end{table}

Os resultados mostram que o maior fluxo foi obtido na posição 1, mais próxima da bobina de alimentação. A posição 2 apresentou fluxo intermediário, enquanto a posição 3 apresentou o menor valor. Esse comportamento indica que o fluxo magnético se distribui pelo núcleo, mas sua intensidade diminui nos ramos mais afastados da excitação.

Para verificar a lei dos nós magnética, considera-se que o fluxo principal no ramo de alimentação deve ser aproximadamente igual à soma dos fluxos nos demais ramos. Assim,

\begin{equation} 
    \Phi_1 \approx \Phi_2 + \Phi_3 . 
    \label{eq:nos_lateral} 
\end{equation}

Substituindo os valores experimentais,

\begin{equation} 
    \Phi_2 + \Phi_3 = 
    2,39 \times 10^{-5} 
    + 
    7,96 \times 10^{-6} 
    = 
    3,19 \times 10^{-5} 
    \; \text{Wb}. 
    \label{eq:soma_lateral} 
\end{equation}

Comparando com

\begin{equation} 
    \Phi_1 = 
    3,71 \times 10^{-5} 
    \; \text{Wb}, 
    \label{eq:phi1_lateral} 
\end{equation}

obtém-se uma diferença relativa de

\begin{equation} 
    \varepsilon = 
    \frac{3,71 \times 10^{-5} - 3,19 \times 10^{-5}} 
    {3,71 \times 10^{-5}} 
    \cdot 100 
    = 
    14,0\%. 
    \label{eq:erro_lateral} 
\end{equation}

A diferença observada é compatível com as limitações do experimento. Parte do fluxo pode ter se dispersado pelo ar, além de existirem erros de posicionamento da bobina de teste, incertezas de leitura dos instrumentos e não idealidades do material magnético.

\subsection{Ensaio E+I com alimentação no ramo central} Na segunda configuração do núcleo E+I, a bobina de alimentação foi posicionada no ramo central. Os resultados são apresentados na Tabela 
\ref{tab:ei_central}.

\begin{table}[H] 
    \centering 
    \caption{Resultados experimentais do núcleo E+I com alimentação no ramo central.} \label{tab:ei_central} 
    \begin{tabular}{ccc} 
        \toprule Posição da bobina de teste & Tensão medida (V) & Fluxo calculado (Wb) \\ 
        \midrule 
        1 & 6 & \(1,59 \times 10^{-5}\) \\ 
        2 & 14 & \(3,71 \times 10^{-5}\) \\ 
        3 & 6 & \(1,59 \times 10^{-5}\) \\ 
        \bottomrule 
    \end{tabular} 
\end{table}

Neste caso, observa-se que os fluxos nos ramos laterais foram aproximadamente iguais. Esse resultado é coerente com a geometria do núcleo, pois, com a bobina de alimentação posicionada no ramo central, o circuito magnético apresenta simetria aproximada entre os caminhos laterais.

A verificação da lei dos nós pode ser feita por

\begin{equation} 
    \Phi_2 \approx \Phi_1 + \Phi_3 . 
    \label{eq:nos_central} 
\end{equation}

Substituindo os valores medidos,


\begin{equation} 
    \Phi_1 + \Phi_3 = 
    1,59 \times 10^{-5}
    + 
    1,59 \times 10^{-5} 
    = 
    3,18 \times 10^{-5} 
    \; \text{Wb}.
    \label{eq:soma_central} 
\end{equation}

Comparando com

\begin{equation} 
    \Phi_2 = 
    3,71 \times 10^{-5}
    \; \text{Wb}, 
    \label{eq:phi2_central} 
\end{equation}

obtém-se

\begin{equation} 
    \varepsilon = 
    \frac{3,71 \times 10^{-5} - 3,18 \times 10^{-5}} 
    {3,71 \times 10^{-5}} 
    \cdot 100 
    =
    14,3\%. 
    \label{eq:erro_central} 
\end{equation}

A diferença percentual foi semelhante à observada no caso com alimentação no ramo lateral. Isso sugere que as principais fontes de erro estão associadas a limitações gerais da montagem experimental, como dispersão de fluxo, folgas entre as peças do núcleo, incertezas de medição e simplificações no modelo teórico.

\subsection{Resultados da simulação FEMM}

A simulação no FEMM foi realizada com o núcleo E+I, utilizando material magnético linear com permeabilidade relativa \(\mu_r = 1000\), conforme indicado na metodologia da videoaula utilizada como apoio. A corrente configurada na bobina foi de \(\SI{0,014}{A}\), compatível com a corrente medida no ensaio experimental com núcleo fechado. A Figura \ref{fig:femm_mapa_fluxo} apresenta o mapa de densidade de fluxo magnético obtido na simulação.

\figuraopcional 
    {femm_mapa_fluxo.jpeg} 
    {0.95} 
    {Distribuição da densidade de fluxo magnético no núcleo E+I obtida no FEMM.} {fig:femm_mapa_fluxo}

A Figura \ref{fig:femm_circuit_properties} apresenta a configuração do circuito na simulação, indicando a corrente total de \(\SI{0,014}{A}\).

\figuraopcional 
    {femm_circuit_properties.jpeg} 
    {0.75} 
    {Propriedades do circuito configurado no FEMM.} 
    {fig:femm_circuit_properties}

A Figura \ref{fig:femm_material} apresenta a configuração do material magnético utilizado na simulação.

\figuraopcional 
    {femm_material.jpeg}
    {0.75} 
    {Material magnético linear utilizado no FEMM com \(\mu_r = 1000\).} 
    {fig:femm_material}

A análise visual do mapa de densidade de fluxo mostra que o maior valor ocorre no ramo onde a bobina de alimentação está posicionada. Nos demais ramos, a densidade de fluxo diminui gradativamente. Esse comportamento está de acordo com os resultados experimentais obtidos para a alimentação no ramo lateral, em que as tensões medidas foram \(\SI{14}{V}\), \(\SI{9}{V}\) e \(\SI{3}{V}\), respectivamente.

Além da análise qualitativa, foram coletados valores pontuais da densidade de fluxo magnético nos três ramos do núcleo. Os valores observados são apresentados na Tabela \ref{tab:femm_b}.

\begin{table}[H] 
    \centering 
    \caption{Valores pontuais de densidade de fluxo magnético obtidos no FEMM.} \label{tab:femm_b} 
    \begin{tabular}{lc} 
        \toprule 
        Ramo analisado & \(|B|\) (T) \\ 
        \midrule 
        Ramo lateral esquerdo & \(2,31443 \times 10^{-4}\) \\ 
        Ramo central & \(1,35510 \times 10^{-4}\) \\ 
        Ramo lateral direito & \(2,22954 \times 10^{-5}\) \\ 
        \bottomrule 
    \end{tabular} 
\end{table}

Normalizando os valores pelo maior valor de densidade de fluxo, obtém-se a Tabela \ref{tab:comparacao_normalizada}.

\begin{table}[H] 
    \centering 
    \caption{Comparação normalizada entre resultados experimentais e FEMM.} \label{tab:comparacao_normalizada} 
    \begin{tabular}{lcc} 
        \toprule 
        Ramo & Resultado experimental normalizado & FEMM normalizado \\ 
        \midrule
        Ramo lateral esquerdo & 1,00 & 1,00 \\
        Ramo central & 0,64 & 0,59 \\ 
        Ramo lateral direito & 0,21 & 0,10 \\ 
        \bottomrule
    \end{tabular} 
\end{table}

A comparação normalizada mostra que o FEMM reproduziu corretamente a tendência observada experimentalmente. Em ambos os casos, o maior fluxo ocorre no ramo lateral esquerdo, seguido pelo ramo central e, por fim, pelo ramo lateral direito. A diferença entre os valores experimentais e simulados pode ser atribuída às simplificações adotadas na simulação, especialmente ao uso de material linear, à aproximação geométrica do núcleo, à ausência de modelagem detalhada dos entreferros e ao fato de os valores de \(B\) terem sido extraídos pontualmente, e não por integração da densidade de fluxo ao longo de toda a seção transversal de cada ramo.

\subsection{Limitações da comparação}

A comparação entre os resultados experimentais e simulados apresenta algumas limitações importantes. Primeiramente, o material magnético real pode apresentar comportamento não linear, histerese e saturação, enquanto a simulação considerou uma curva B-H linear com \(\mu_r = 1000\). Além disso, o contato entre as peças do núcleo não é ideal. Pequenos entreferros podem surgir devido a folgas mecânicas, rugosidade das superfícies e desalinhamentos na montagem.

Outra limitação relevante é que a bobina de teste foi utilizada para medir indiretamente o fluxo magnético por meio da tensão induzida. Essa medição depende do correto posicionamento da bobina, da leitura do voltímetro e da hipótese de operação senoidal. No FEMM, por outro lado, os valores de densidade de fluxo foram extraídos em pontos específicos do núcleo, o que não representa necessariamente o valor médio em toda a seção transversal.

Apesar dessas limitações, os resultados obtidos foram coerentes. O experimento mostrou que o fechamento do núcleo reduz a relutância do circuito magnético, aumentando o fluxo e reduzindo a corrente de magnetização. No núcleo E+I, observou-se a divisão do fluxo entre os ramos, com comportamento aproximadamente compatível com a lei dos nós magnética. A simulação FEMM confirmou qualitativamente essa distribuição.

% CONCLUSÃO%

\section{Conclusões}

O experimento permitiu verificar a influência da relutância no comportamento de circuitos magnéticos. No ensaio com núcleo U+I, observou-se que o fechamento do caminho magnético aumentou a tensão induzida na bobina de teste de aproximadamente \(\SI{4}{V}\) para \(\SI{12}{V}\). Como a tensão induzida é proporcional ao fluxo magnético, conclui-se que o fluxo no núcleo fechado foi cerca de três vezes maior que no núcleo aberto.

Também foi observada uma redução significativa da corrente de alimentação, que passou de \(\SI{0,120}{A}\) no núcleo aberto para \(\SI{0,014}{A}\) no núcleo fechado. Esse resultado evidencia que a redução da relutância permite produzir maior fluxo com menor corrente de magnetização.

No ensaio com núcleo E+I, verificou-se que o fluxo magnético se distribui entre os ramos do núcleo. Quando a bobina de alimentação foi posicionada no ramo lateral, o maior fluxo ocorreu no ramo mais próximo da excitação. Quando a bobina foi posicionada no ramo central, os fluxos nos ramos laterais apresentaram valores aproximadamente iguais, indicando uma divisão simétrica do fluxo.

A lei dos nós magnética foi verificada de forma aproximada. Em ambos os casos analisados, a diferença entre o fluxo principal e a soma dos fluxos nos demais ramos foi próxima de 14\%. Essa diferença é aceitável considerando as limitações experimentais, como dispersão de fluxo, entreferros não intencionais, incertezas de medição e não idealidades do material.

A simulação realizada no FEMM apresentou distribuição de fluxo coerente com os resultados experimentais. O mapa de densidade de fluxo mostrou maior concentração no ramo onde estava localizada a bobina de alimentação e redução gradual nos demais ramos. Embora a simulação tenha utilizado um modelo simplificado, com material linear e permeabilidade relativa constante, ela foi útil para visualizar o comportamento do circuito magnético e reforçar a interpretação dos resultados experimentais.


\begin{thebibliography}{99} 
    \bibitem{chapman2013} 
    CHAPMAN, Stephen J. 
    \textit{Fundamentos de Máquinas Elétricas}. 
    5. ed. 
    Porto Alegre: AMGH Editora, 2013. 
    
    \bibitem{umans2014} 
    UMANS, Stephen D. 
    \textit{Máquinas Elétricas de Fitzgerald e Kingsley}. 
    7. ed. 
    Porto Alegre: AMGH Editora, 2014.
    \bibitem{femm} MEEKER, David. 
    \textit{Finite Element Method Magnetics: FEMM}. 
    Software de análise por elementos finitos aplicado a problemas magnéticos. 
    
\end{thebibliography} 
    
\end{document}